\scanpage{45}
\mspage{- 23 -}
% archive sheet number: 44

\[
Du_i(x;z_1,..,z_q)=f_i(x;u_1(x;z_1,..,z_q),..,u_p(x;z_1,..,z_q);z_1,..,z_q)\tag*{(12)}
\]
(D \uline{dérivation par rapport à} x, $1\handwritten{\leq}i\handwritten{\leq}p$) \uline{et}
\[
u_i(a;z_1,..,z_q)=b_i\qquad(1\handwritten{\leq}i\handwritten{\leq}p)\quad\text{\uline{pour tout}}\quad(z_1,..,z_q)\handwritten{\in}W\tag*{(13)}
\]

Par translation, on peut se ramener au cas où $a=b_j=c_k=0$ pour $1\handwritten{\leq}j\handwritten{\leq}p$ et $1\handwritten{\leq}k\handwritten{\leq}q$ ; chacune des fonctions $f_i$ est alors égale à la somme d'une série entière sommable dans un polycylindre ouvert $B\handwritten{\times}P\handwritten{\times}Q$, où $B\handwritten{\subset}K$, $P\handwritten{\subset}K^p$ et $Q\handwritten{\subset}K^q$
\[
\begin{aligned}
f_i(x;y_1,..,y_p;z_1,..,z_q)&=\handwritten{\sum}g_i(n;r_1,..,r_p;\handwritten{\lambda}_1,..,\handwritten{\lambda}_q)\mathbin{.}\\
&\quad\mathbin{.}x^n y_1^{r_1}..y_p^{r_p}z_1^{\handwritten{\lambda}_1}z_2^{\handwritten{\lambda}_2}\ldots z_q^{\handwritten{\lambda}_q}.
\end{aligned}\tag*{(14)}
\]

Pour établir l'existence des fonctions $u_i$, nous devons prouver qu'il existe p séries entières sans terme indépendant de x
\[
u_i(x;z_1,..,z_q)=\handwritten{\sum}h_i(m;\handwritten{\mu}_1,..,\handwritten{\mu}_q)x^m z_1^{\handwritten{\mu}_1}\ldots z_q^{\handwritten{\mu}_q}
\]
sommables dans un polydylindre $U\handwritten{\times}W$ de $K^{q+1}$, substituables dans les séries entières (14) et satisfaisant à (12). Si le problème est possible, les séries entières par rapport à x et aux $z_{\handwritten{\ell}}$ $(1\handwritten{\leq\ell\leq}q)$ obtenues par cette substitution dans les deux membres de (12) doivent avoir mêmes coefficients, ce qui donne les conditions
\[
\left.\begin{aligned}
h_i(0;\handwritten{\mu}_1,..,\handwritten{\mu}_q)&=0\\
mh_i(m;\handwritten{\mu}_1,\ldots,\handwritten{\mu}_q)&=\handwritten{\sum_{n,(r_j),(\lambda_\ell)}}\unsure{g|G}_i(n;r_1,..,r_p;\handwritten{\lambda}_1,\ldots,\handwritten{\lambda}_q)\mathbin{.}\\
&\quad\mathbin{.}\handwritten{\prod_{j=1}^p\prod_{k=1}^{r_j}}h_i(s_{kj};\handwritten{\nu}_{kj1},\ldots,\handwritten{\nu}_{kjq})
\end{aligned}\right\}\tag*{(15)}
\]
pour tout entier $m\handwritten{>}0$ et tout système d'entiers $\handwritten{\mu_\ell\geq}0$ $(1\handwritten{\leq\ell\leq}q)$ ; la sommation du second membre des relations (15) doit être étendue à tous les systèmes d'entiers $\handwritten{\geq}0$, $n,r_j,\handwritten{\lambda_\ell}$, $s_{kj}$ et $\handwritten{\nu}_{kj\handwritten{\ell}}$ satisfaisa\handwritten{-}nt aux conditions
